% ============================================================
% Meta-Labeling Two-Stage Architecture Flowchart
% CSC491 Textbook, Chapter: Machine Learning
%
% Preamble requirements:
%     \usepackage{tikz}
%     \usetikzlibrary{arrows.meta, positioning, fit, backgrounds}
%     \usepackage{xcolor}
%
%     \definecolor{pipenavy}{RGB}{26,  47,  89}
%     \definecolor{pipeteal}{RGB}{17, 150, 150}
%     \definecolor{pipegreen}{RGB}{31, 120,  85}
%     \definecolor{pipegold}{RGB}{160,110,   0}
%     \definecolor{pipegray}{RGB}{100,100, 100}
% ============================================================

\begin{figure}[ht]
\centering
\begin{tikzpicture}[
  node distance   = 0.7cm and 1.1cm,
  %% --- Box styles ---
  databox/.style = {
    rectangle, rounded corners=3pt,
    minimum width=2.6cm, minimum height=1.1cm,
    align=center, text=white, font=\sffamily\small,
  },
  labelbox/.style = {
    rectangle, rounded corners=3pt,
    minimum width=2.0cm, minimum height=0.85cm,
    align=center, font=\sffamily\scriptsize,
    draw, line width=0.6pt,
  },
  stagebox/.style = {
    rectangle, rounded corners=4pt,
    draw=black!25, line width=0.5pt,
    inner sep=8pt, dashed,
  },
  %% --- Arrow style ---
  arr/.style={
    ->, -{Stealth[length=6pt, width=4pt]},
    line width=1.4pt, color=black!40,
  },
]

%% ── Row 1: main pipeline nodes ──────────────────────────────

\node[databox, fill=pipenavy] (feat)
  {\textbf{Feature}\\\textbf{Matrix}\\[-2pt]\scriptsize\color{white!70!black}RSI, ATR, EMA \ldots};

\node[databox, fill=pipeteal, right=of feat] (primary)
  {\textbf{Primary}\\\textbf{Model}\\[-2pt]\scriptsize\color{white!70!black}RSI crossover rule};

\node[databox, fill=pipegold!80!black, right=of primary] (side)
  {\textbf{Trade Side}\\[-2pt]\scriptsize\color{white!85!black}Long ($+1$)\\[-1pt]\scriptsize\color{white!85!black}Short ($-1$)};

\node[databox, fill=pipegreen, right=of side] (secondary)
  {\textbf{Secondary}\\\textbf{Model}\\[-2pt]\scriptsize\color{white!70!black}Decision tree};

\node[databox, fill=pipenavy!80, right=of secondary] (metalabel)
  {\textbf{Meta-Label}\\[-2pt]\scriptsize\color{white!85!black}Take ($1$)\\[-1pt]\scriptsize\color{white!85!black}Pass ($0$)};

\node[databox, fill=pipegreen!70!black, right=of metalabel] (size)
  {\textbf{Position}\\\textbf{Sizing}\\[-2pt]\scriptsize\color{white!70!black}Full / reduced / skip};

%% ── Arrows between main nodes ───────────────────────────────

\draw[arr] (feat)      -- (primary);
\draw[arr] (primary)   -- (side);
\draw[arr] (side)      -- (secondary);
\draw[arr] (secondary) -- (metalabel);
\draw[arr] (metalabel) -- (size);

%% ── Stage labels below each node ────────────────────────────

\node[font=\sffamily\scriptsize, text=pipeteal,   below=0.30cm of primary]  {Stage 1: Direction};
\node[font=\sffamily\scriptsize, text=pipegreen,  below=0.30cm of secondary]{Stage 2: Conviction};
\node[font=\sffamily\scriptsize, text=pipegold!80!black, below=0.30cm of side]{Signal};
\node[font=\sffamily\scriptsize, text=pipenavy!80,below=0.30cm of metalabel]{Filter};
\node[font=\sffamily\scriptsize, text=pipegreen!70!black, below=0.30cm of size]{Output};

%% ── Feature matrix also feeds secondary model (dashed arc) ──
%% Shows that the same features inform both models

\draw[
  ->, -{Stealth[length=5pt, width=3.5pt]},
  dashed, line width=1.0pt, color=pipenavy!50,
  bend right=38,
] (feat.south) to node[below, font=\sffamily\scriptsize, text=pipenavy!60,
    yshift=-2pt] {same features inform both models}
  (secondary.south);

\end{tikzpicture}

\caption{%
  The Meta-Labeling architecture. The \textbf{Primary Model} (Stage~1) uses a
  rule-based signal to determine the \textit{direction} of the trade (long or short).
  The \textbf{Secondary Model} (Stage~2) then acts as a filter, predicting whether the primary signal
  is likely to be correct.
  The dashed arc highlights that both models draw from the same feature matrix,
  but answer fundamentally different questions.%
}
\label{fig:meta-labeling-flow}
\end{figure}